% Related literature for the Nera draft (Oct 23, 2026). Base: the author's text in
% latex/JMP_DS_draft/sections/01_introduction.tex (strands two to four keep his sentences).
% Added and checked on the PDFs: Kearney, Levine and Pardue (2022) pp. 151-153 and the section on shifting
% priorities; Fernandez-Villaverde and Norrick (2026) abstract; Doepke et al. (2023) contents and Section 5 title;
% Coven et al. (2025) abstract; Davis and Van Nieuwerburgh (2015) and Piazzesi and Schneider (2016) abstracts;
% Chambers, Garriga and Schlagenhauf (2009) abstract; Sommer and Sullivan (2018) abstract; Saiz (2010) and
% Glaeser and Gyourko (2018) first pages.
\paragraph{Related literature.}
This paper sits at the intersection of several branches of the literature.

First, it relates to work on the recent fall in fertility. In the United
States, the birth rate fell by about 20\% between 2007 and 2020.
\citet{kearneyPuzzleFallingUS2022} find no strong link between this decline
and specific policies or economic conditions beyond the Great Recession, and
point instead to shifting priorities across cohorts.
\citet{fernandezVillaverdeNorrickTerraIncognita2026} document that the world
as a whole is likely below replacement and argue that nothing in an economy
pushes fertility back to 2.1. In high-income countries, the old negative
relationships between income, women's work and fertility have weakened, and
the ease of combining career and family has become a central determinant of
fertility \citep{doepkeEconomicsFertilityNew2023}. I take the decline as given
and ask a different question: how much housing policy can raise births from
here on.

Second, the paper relates to the literature on housing costs and fertility. A
large empirical literature finds that housing conditions shape childbearing.
Fertility is lower where rent per room is higher
\citep{simonTamuraHigherRents2009}, dwelling size predicts first births
\citep{stromHousingFirstBirths2010}, and higher house prices raise births
among owners but lower them among renters
\citep{dettlingHousePricesBirth2014,lovenheimFamilyWealthShocks2013}. Access
to mortgage credit also raises births
\citep{hacamoBabiesMortgageDeregulation2021,cummingMonetaryPolicyBirth2024,dettlingKearneyModernMortgage2025,vandoornikFazioRamadoraiSkrastinsHousingFertility2025}.
Two recent papers model housing and fertility jointly.
\citet{khorunzhinaMillerAmericanDream2022} estimate a dynamic model of female
labor supply, fertility and the first home purchase at given house prices,
and \citet{couillardBuildBabyBuild2025} estimates a model of fertility,
living arrangements and location in which rents are the cost of housing.
Unlike these papers, I study fertility together with saving, tenure and the
choice of space under uninsurable earnings risk, in a housing market where
the price is determined in equilibrium. This allows me to separate the
recurring cost of family space from the equity needed to buy it, and to
compare housing supply and mortgage credit within the same economy.

Third, the paper relates to work on earnings risk and fertility. Because a
child is an irreversible commitment, prudent parents who face more risk want
fewer children \citep{fraserIncomeRisk2001}. Quantitative life-cycle models
find that uninsurable earnings risk delays and reduces births
\citep{sommerFertilityChoiceLife2016,ejrnaesJorgensenFamilyPlanning2020,nakakuniFertilityIncomeUncertainty2026},
and \citet{deLaCroixPommeretChildbearing2021} show that risk gives the
postponement of a first birth an option value. Quasi-experimental studies
find that job displacement and job insecurity lower fertility
\citep{delBonoWeberWinterEbmer2012,clarkLepinteurJobInsecurity2022}. None of
these papers includes housing. In my model, part of the commitment a child
entails is the space it needs, paid for every period at a price set in the
housing market. The response of fertility to earnings risk therefore depends
on housing costs and on the down payment required to own family space.

Fourth, the paper builds on quantitative models of housing in macroeconomics,
reviewed by \citet{davisHousingFinanceMacroeconomy2015} and
\citet{piazzesiSchneiderHousing2016}. These models study tenure choice under
earnings risk, down payments and transaction costs
\citep{fisherGervaisHomeOwnership2011,halketVasudevSaving2014,boarGoreaMidriganLiquidity2022},
and the joint determination of house prices, rents and ownership
\citep{kaplanMitmanViolanteHousing2020,greenwaldGurenCredit2025,greaneyDynamicUrbanEconomics2025}.\footnote{In
\citet{chambersAccountingChangesHomeownership2009}, mortgage innovation
accounts for most of the rise in homeownership between 1994 and 2005, and in
\citet{sommerImplicationsUSTax2018}, eliminating the mortgage interest
deduction lowers house prices and raises homeownership. On the supply side,
\citet{saizGeographicDeterminantsHousing2010} relates housing supply to land
scarcity across metropolitan areas, \citet{glaeserEconomicImplicationsHousing2018}
compare house prices with construction costs, and the elasticity I use is the
floor-space estimate of \citet{baum-snowMicrogeographyHousingSupply2024}.} In
this literature, rising earnings risk accounts for much of the decline in
homeownership among young households
\citep{fisherGervaisHomeOwnership2011,pazPardoHomeownership2024}. Closest to my
property-tax experiment, \citet{coven2024} show that higher property taxes
lower the down payment through capitalization and raise homeownership among
young households. In my model, a rebated property tax also lowers the price,
but it raises the per-period cost of space, and births barely move. Household
size, however, is exogenous or absent in these models. I make fertility a
choice. A birth raises the household's demand for space over its life cycle
and, a generation later, the number of households that compete for housing.
Housing policy therefore shapes both family size and the future demand for
housing.
